\section{Smallest measurable force of translational fixture}
\label{sec:smallestForce}

For a specimen of $m = 0.1\,\text{g} = 1\times10^{-4}\,\text{kg}$, the smallest measurable magnetic force is given by

\begin{equation}
	F_{m,\min} = m g \frac{z}{L}
\end{equation}

\noindent where $g = 9.81\,\text{m/s}^2$, $z = 1\times10^{-3}\,\text{m}$, and $L = 1\,\text{m}$. Substituting these values yields

\begin{align}
	F_{m,\min} &= (1\times10^{-4})(9.81)\left(\frac{1\times10^{-3}}{1}\right) \nonumber\\
	&= 9.81\times10^{-7}\,\text{N} \approx 1.0\times10^{-6}\,\text{N}.
\end{align}

\noindent Thus, for a $0.1\,\text{g}$ mass, the smallest detectable magnetic force is approximately $1\,\mu\text{N}$. The corresponding force ratio remains 
\begin{equation}
	\frac{F_m}{F_g} = \frac{z}{L} = 10^{-3},
\end{equation}
\noindent indicating that the setup can detect magnetic forces as small as one-thousandth of the object’s weight.
